% ==========================================
% TABEL VI.1 HASIL FUNCTIONALITY TESTING
% ==========================================
\begin{longtable}{|p{1.2cm}|p{4.5cm}|p{5cm}|p{2cm}|}
\caption{Hasil \textit{Functionality Testing}} \label{tab:functionality-testing} \\

\hline
\textbf{ID} & \textbf{Skenario Pengujian} & \textbf{Hasil yang Diharapkan} & \textbf{Status} \\ \hline
\endfirsthead

\caption[]{Hasil \textit{Functionality Testing} (lanjutan)} \\
\hline
\textbf{ID} & \textbf{Skenario Pengujian} & \textbf{Hasil yang Diharapkan} & \textbf{Status} \\ \hline
\endhead

\hline
\multicolumn{4}{r}{\textit{Bersambung ke halaman berikutnya}} \\
\endfoot

\hline
\endlastfoot

\multicolumn{4}{|c|}{\textbf{Mobile (Pelari)}} \\ \hline

TC-01 & Pelari bergerak di sepanjang rute & ETA, jarak tersisa, waktu tempuh, dan \textit{pace} terhitung serta diperbarui secara dinamis & Berhasil \\ \hline

TC-02 & Pelari menekan tombol pusatkan kamera (\textit{center}) & Peta otomatis bergeser dan memusatkan tampilan ke titik lokasi pelari saat ini & Berhasil \\ \hline

TC-03 & Pelari menekan tombol \textit{emergency} (SOS) & Tampilan darurat muncul dan hitung mundur peringatan dimulai & Berhasil \\ \hline

TC-04 & Pelari menekan tombol menu utama & Laci menu terbuka dan menampilkan opsi navigasi tambahan & Berhasil \\ \hline

TC-05 & Pelari menekan tombol info tim & Panel informasi muncul menampilkan data tim dan anggota & Berhasil \\ \hline

TC-06 & Pelari mengubah tuas (\textit{switch}) \textit{visibility} & Status visibilitas berubah antara \textit{public} atau \textit{private} & Berhasil \\ \hline

TC-07 & Pelari menekan tombol \textit{logout} & Sesi pengguna berakhir dan aplikasi kembali ke layar \textit{login} & Berhasil \\ \hline

TC-08 & Pelari menekan layar (\textit{tap}) 2x saat kondisi \textit{emergency} & Kondisi \textit{emergency} dibatalkan (\textit{cancel}) dan kembali ke mode peta & Berhasil \\ \hline

TC-09 & Pelari menekan navigasi arah (sebelum/selanjutnya) & Peta berpindah dan memfokuskan (\textit{spectating}) pelari lain dalam tim yang sama & Berhasil \\ \hline


\multicolumn{4}{|c|}{\textbf{Web Publik (Spectator)}} \\ \hline

TC-10 & Pengguna membuka halaman web publik & Peta menampilkan titik pelari-pelari yang membuka visibilitas publik & Berhasil \\ \hline

TC-11 & Pengguna mengeklik salah satu titik pelari & Muncul detail informasi pelari beserta nama timnya & Berhasil \\ \hline

TC-12 & Pengguna mencari tim dari daftar direktori tim & Peta otomatis bergeser mengarahkan tampilan ke titik tim tersebut & Berhasil \\ \hline

TC-13 & Pengguna mencari \textit{water station} dari direktori & Peta otomatis bergeser mengarahkan tampilan ke titik lokasi \textit{water station} & Berhasil \\ \hline


\multicolumn{4}{|c|}{\textbf{Web Privat (Ketua Tim)}} \\ \hline

TC-14 & Sistem mendeteksi kondisi \textit{emergency} atau tersesat & Antarmuka Ketua menerima notifikasi masuk & Berhasil \\ \hline

TC-15 & Ketua mengeklik titik pelari pada peta & Peta menampilkan posisi pelari dan rincian detailnya & Berhasil \\ \hline

TC-16 & Ketua mengakses opsi bagikan & Tautan untuk \textit{supporter} dibuat dan siap disalin & Berhasil \\ \hline


\multicolumn{4}{|c|}{\textbf{Web Privat (Supporter)}} \\ \hline

TC-17 & Sistem mendeteksi kondisi \textit{emergency} atau tersesat & Antarmuka Supporter menerima notifikasi peringatan & Berhasil \\ \hline

TC-18 & Supporter mengeklik titik pelari yang dipantau & Peta menampilkan posisi pelari dan rincian detailnya & Berhasil \\ \hline


\multicolumn{4}{|c|}{\textbf{Web Privat (Panitia)}} \\ \hline

TC-19 & Panitia mengeklik titik pelari pada peta & Peta menampilkan posisi detail pelari dari seluruh tim & Berhasil \\ \hline

TC-20 & Panitia membuka daftar informasi tim & Menampilkan rincian lengkap dari tim dan anggotanya & Berhasil \\ \hline

TC-21 & Panitia membuka modul daftar pelari & Menampilkan daftar status dari masing-masing pelari & Berhasil \\ \hline

TC-22 & Panitia membuka dasbor statistik & Tampil rekapitulasi angka statistik perlombaan & Berhasil \\ \hline

TC-23 & Sistem mendeteksi adanya pelari yang SOS/tersesat & Peringatan darurat masuk ke halaman panitia dengan penanda merah & Berhasil \\ \hline

\end{longtable}
